% Generated draft: requires mathematical and source review.
\documentclass[11pt]{article}
\usepackage[margin=25mm]{geometry}
\usepackage{fontspec}
\setmainfont{Latin Modern Roman}
\usepackage{amsmath,amssymb,mathtools,unicode-math}
\setmathfont{Latin Modern Math}
\usepackage{enumitem,needspace,xurl,hyperref}
\setlength{\parindent}{0pt}
\setlength{\parskip}{5pt}
\setlength{\emergencystretch}{3em}
\newcommand{\N}{\mathbb{N}}
\newcommand{\Z}{\mathbb{Z}}
\newcommand{\Q}{\mathbb{Q}}
\newcommand{\R}{\mathbb{R}}
\newcommand{\C}{\mathbb{C}}
\newcommand{\F}{\mathbb{F}}
\newcommand{\A}{\mathbb{A}}
\newcommand{\PP}{\mathbb{P}}
\newcommand{\E}{\mathbb{E}}
\newcommand{\eps}{\varepsilon}
\newcommand{\dd}{\,\mathrm{d}}
\begin{document}
\section*{UnsolvedMath ID: 20001298}
\textbf{Source problem code:} catalogue entry
\subsection{Definitions and mathematical statement}
For \(n\geq1\), use the standard Gaussian density
\[
 \varphi_n(x)=(2\pi)^{-n/2}e^{-|x|^2/2}.
\]
For a full-dimensional bounded ellipsoid
\[
 E=\{x\in\R^n:(x-a)^{\mathsf T}A^{-1}(x-a)\leq1\},
 \quad a\in\R^n,\quad A=A^{\mathsf T}>0,
\]
define its Gaussian surface area, or Gaussian perimeter, by
\[
 P_\gamma(E)=\int_{\partial E}\varphi_n(x)\,\dd\mathcal H^{n-1}(x).
\]
In dimension one, \(\mathcal H^0\) counts the two endpoints. Equivalently for these smooth bodies this is the outer Gaussian Minkowski-content derivative. The source question is whether there is an absolute constant \(C\), independent of \(n,a,A\), such that \(P_\gamma(E)\leq C\) for every ellipsoid.

The supplied aid reports the sharp affirmative statement
\[
 P_\gamma(E)<\sqrt{2/\pi}\quad\text{for every full-dimensional bounded ellipsoid},
 \qquad
 \sup_{E\subseteq\R^n\text{ such ellipsoids}}P_\gamma(E)=\sqrt{2/\pi}
 \quad(n\geq1).
\]
Translated ellipsoids are included. No ellipsoid in this class attains the supremum. The density normalization is material, and the optimization is already sharp in each fixed dimension.

The reported limiting examples in dimension one are intervals \([-r,r]\) as \(r\downarrow0\). For fixed \(n\geq2\), they are ellipsoids
\[
 E_{s,R}=\{x:x_1^2/s^2+(x_2^2+\cdots+x_n^2)/R^2\leq1\},\qquad s,R>0.
\]
For fixed \(s\), the aid reports convergence of their perimeters as \(R\to\infty\) to the perimeter \(2\varphi_1(s)\) of the slab \(|x_1|\leq s\); then \(s\downarrow0\) gives \(\sqrt{2/\pi}\). These are ordered limits of bounded nondegenerate ellipsoids, not an attained maximum at a slab or a lower-dimensional set.

The original remark contrasts the general-convex-body bound \(O(n^{1/4})\), attributed to Keith Ball, and says that intersections of random halfspaces attain that order. The aid attributes the historical affirmative ellipsoid solution to Kane's degree-\(d\) polynomial-threshold theorem and reports a specialized Gaussian Crofton argument for the quadratic case and strict nonattainment. These are supplied attributions, not newly verified literature claims.
\subsection{Short English statement}
The aid reports that every bounded nondegenerate ellipsoid, including translated ones, has standard Gaussian perimeter strictly below \(\sqrt{2/\pi}\). That constant is the supremum in every fixed dimension, approached through degenerating limits but never attained by an ellipsoid in the stated class.
\subsection{Sources}
\begin{itemize}
\item Local registry: permanent UnsolvedMath ID \(20001298\), exact public code \texttt{catalogue entry}; \texttt{aim-convex-geometry-notes.json}, \texttt{notes[29]}.
\item AIM workshop ``Algorithmic convex geometry'', Problem 11, attributed to Ryan O'Donnell: \url{https://aimath.org/WWN/convexgeometry/convexgeometry.pdf}. The original remark attributes the general convex-body \(O(n^{1/4})\) bound to Keith Ball and records random-halfspace examples of that order. The final 3 is a page artifact. The literature field is empty.
\item Research attempt 1: \texttt{full\_solution}, status at run \texttt{solved}, rigor label \texttt{proved}, novelty confidence low, classification \texttt{SOLVED-IN-LITERATURE}. It names Kane's polynomial-threshold theorem and a specialized Gaussian Crofton proof, and reports the fixed-dimension sharp supremum with strict nonattainment as a candidate sharp-bound contribution. No complete Kane, Ball or Crofton bibliography is supplied. The current registry status is solved.
\end{itemize}
\end{document}
